\section{Part B: the lifting assertion D-031}
\label{sec:B-audit}

\subsection{The source statement and conventions}

We use cochain complexes: $d:X^i\to X^{i+1}$, and
$X[r]^i=X^{i+r}$ with differential $(-1)^r d_X$.
All tensor products and linear maps below are over a fixed commutative
unital ring $k$.  No flatness over $\mathbb Z$ is assumed.
The paper's coefficient ring $A$ is our $k$; its algebra $A^*$ is our $R$.
A PD dg algebra is a triple $(R,J,\gamma)$: $R$ is strictly graded
commutative, $J$ is a differential ideal, and the divided powers on
$J_{\rm even}$ satisfy the PD identities, including
$\gamma_n(xy)=0$ for odd homogeneous $x\in R$, $y\in J$, $n>1$, and
$d\gamma_n(x)=\gamma_{n-1}(x)dx$.

The following locators refer to the downloaded version-4 PDF
\cite{B_CaradotLin2026}.  Its printed numbering controls our citations.
The HTML assigns the three displayed equations below the numbers
$(14),(15),(16)$ instead of the PDF's $(11),(12),(13)$.

\begin{center}
\begin{tabular}{lll}
\hline
Material & Printed page & Source line(s)\\
\hline
\S2.1, cochain convention and termwise freeness & 6 & 354--369\\
\S2.2, eight PD identities & 6--7 & 373--389\\
PBW formula (11) & 16 & 954--956\\
Differential formula (12) & 17 & 961--964\\
Theorem 3.9; extension formula (13) & 19 & 1076--1088\\
Proposition 3.10; its proof & 19--20 & 1111--1143\\
Theorem 3.15 & 23--25 & 1360--1452\\
\hline
\end{tabular}
\end{center}

The freeness definition is ``$M^i$ is a free $A$-module for all $i$''.
Here is the complete quantified mathematical content of Proposition 3.10,
with the above renaming:
\[
\begin{gathered}
(R',J',\gamma'),(R,J,\gamma)\in\mathrm{PDdgAlg}(k),\qquad
\pi:(R',J',\gamma')\longrightarrow(R,J,\gamma),\quad
\pi(J'^i)=J^i\quad(i\in\mathbb Z),\\
X\in\mathrm{Cplx}(k),\quad X^i\text{ free for every }i,
\qquad \rho\in\operatorname{Hom}_{\mathrm{Cplx}(k)}(X,J),\\
\exists\,\widetilde\rho':\mathrm{TS}_k(X)\longrightarrow R'
\text{ in }\mathrm{PDdgAlg}(k),\qquad
\pi\widetilde\rho'=\widetilde\rho.
\end{gathered}                                                    \tag{B.1}
\]
Here $\widetilde\rho$ denotes the extension in Theorem 3.9.
The proposition and its diagram have no equation number.

\begin{proposition}[B-001 / D-031: Refuted]
\label{B001}
The assertion \textup{(B.1)} is false over every nonzero commutative ring
$k$, already for complexes supported in degrees $-2,-1$ and for
square-zero PD ideals.  In particular the refutation D-031 in SA08 is
correct.
\end{proposition}

\begin{proof}
Let $E^{-2}=ku$, $E^{-1}=kv$, $d_Eu=v$, with all other terms zero.
Let $M^{-2}=k\varepsilon$, with all other terms and all differentials zero.
Put
\[
R'=k\oplus E,\quad J'=E,\qquad R=k\oplus M,\quad J=M.
\]
The displayed ideals have zero products.  More explicitly,
$(a,e)(b,f)=(ab,af+be)$; the base summand has degree zero and differential
zero.  This multiplication is associative and unital: both bracketings of
a triple give $(abc,abg+acf+bce)$.  The Koszul commutativity identity holds
because a product of two ideal elements is zero and scalars have degree
zero.  Every odd square vanishes.  The differential is a derivation,
since both sides of its Leibniz identity vanish on two ideal inputs and
the identity on a scalar and an ideal input is $d(ae)=a\,de$.
Also $d^2=0$ and $J',J$ are differential ideals.

Set $\gamma_0(x)=1$, $\gamma_1(x)=x$ and $\gamma_n(x)=0$ for $n\geq2$
on each even ideal.  We check all PD identities.  The product identity
has both sides zero for two positive indices, including $(1,1)$ because
$x^2=0$; index zero gives the unit identity.  Scalar compatibility is
linear for index one and has both sides zero for larger indices.
For the sum identity, every mixed term has two ideal factors and vanishes;
the terms with an index zero give the required identity.  For iterated
divided powers, indices $p=q=1$ give $x$, an index one reduces to the
definition, and $p,q\geq2$ give zero on both sides.  The degree condition
holds because every higher divided power is zero.  Products of odd
elements are zero.  Finally, differential compatibility is the identity
$dx=dx$ for $n=1$; for $n=2$ its right side is $x\,dx=0$; for $n>2$
the right side also vanishes.  These verifications use no division and
apply in characteristics two and three and with arbitrary torsion.

Define a morphism of PD dg $k$-algebras by
\[
\pi:R'\longrightarrow R,\qquad
\pi(1)=1,\quad\pi(u)=\varepsilon,\quad\pi(v)=0.
                                                               \tag{B.2}
\]
It preserves products, the distinguished ideals and every divided power.
It is a chain map because $d_R\pi(u)=0=\pi(v)=\pi d_{R'}u$.
Its restriction $J'\to J$ is surjective in every degree.
Take $X=M$ and $\rho=\operatorname{id}_M:X\to J$.
Each term of $X$ is free, including the zero terms.
A degree-zero linear lift must send $\varepsilon$ to $cu$ with $c=1$.
The chain-map equation gives $cv=0$, hence $c=0$, since $v$ is a basis
element over $k$.  Thus a lift would give $1=0$ in the nonzero ring $k$.

For completeness the algebra map appearing in the conclusion exists
without any use of the disputed general freeness theorem.  The complex
$X$ has one even generator, and
$\mathrm{TS}_k(X)=\bigoplus_{n\geq0}k\varepsilon^{[n]}$, with
$\varepsilon^{[a]}\varepsilon^{[b]}=\binom{a+b}{a}\varepsilon^{[a+b]}$,
zero differential and its usual divided powers.
Send $\varepsilon^{[0]}$ to $1$, $\varepsilon^{[1]}$ to $\varepsilon$,
and every higher basis element to zero.  The preceding PD checks show
that this is $\widetilde\rho:\mathrm{TS}_k(X)\to R$.
Any purported algebra lift restricts on $X$ to the impossible chain lift.
Assigning weight one to $u,v,\varepsilon$ and weight zero to $k$ makes
every map weight preserving; no completion occurs.

The precise failing step in the source proof is the passage from
``As $X^*$ is assumed to be free'' to the existence of a chain lift.
The underlying graded lift here exists uniquely, and its nonzero
differential defect is $\varepsilon\mapsto v$.
\end{proof}

The independent program \texttt{B\_lifting\_certificate.py} verifies
the integer certificate
\[
(1,-1)\binom{1}{1}=0,\qquad (1,-1)\binom{1}{0}=1
\]
for the two lift equations.  It also exhausts the finite algebras over
$\mathbb Z/2$, $\mathbb Z/3$, $\mathbb Z/4$, including multiplication,
the derivation and projection identities, and PD indices through six.
Those finite PD checks supplement the all-index proof above.

\subsection{The free algebra used in the correction}

We write $F_k(X)$ for the free \emph{strict} PD dg algebra on a
termwise-free cochain complex $X$, with its augmentation ideal.
We specify the comparison with the source's symmetric invariants, so
that every use of freeness below has a defined map.

\begin{proposition}[B-002: Proved; explicit free-algebra comparison]
\label{B002}
For a homogeneous basis $B$ of $X$, the underlying graded algebra is
\[
F_k(X)=\Gamma_k(X^{\rm even})\otimes_k\bigwedge_k(X^{\rm odd}).
                                                               \tag{B.3}
\]
Its basis consists of the finite ordered monomials
$x^{[f]}=\prod_{b\in B}b^{[f(b)]}$, with $f(b)\in\mathbb N$ for even
$b$ and $f(b)\in\{0,1\}$ for odd $b$.  The differential extends $d_X$ by
the Leibniz rule and $d(b^{[n]})=b^{[n-1]}d_Xb$ for even $b$.
For every PD dg triple $(R,J,\gamma)$, restriction along
$\iota:X\to F_k(X)^+$ gives a bijection
\[
\operatorname{Hom}_{\mathrm{PDdgAlg}(k)}(F_k(X),(R,J,\gamma))
\;\longrightarrow\;
\operatorname{Hom}_{\mathrm{Cplx}(k)}(X,J).                       \tag{B.4}
\]

There is a natural injective morphism of PD dg algebras
$j_X:F_k(X)\to\mathrm{TS}_k(X)$ carrying every generator to the same
tensor of length one.  With $k[2]=\{a\in k:2a=0\}$, its cokernel as a
graded $k$-module is a direct sum of copies of $k[2]$, one for each
finite multiset of basis elements with a repeated odd basis element.
Thus $j_X$ is an isomorphism for $k[2]=0$, including $k=\mathbb Z$.
\end{proposition}

\begin{proof}
The divided power polynomial construction and its universal property
are those of \cite[Lemma 23.5.1, Tag 07H5]{B_StacksPD}.
An odd variable is adjoined as an exterior variable; its divided powers
on an even expression $a+be$ are
$\gamma_n(a)+\gamma_{n-1}(a)be$, as in
\cite[Example 23.6.2]{B_StacksTate}.
Successive adjunction and the union over finite subsets of $B$ give
\textup{(B.3)}.  The actual integer degrees are assigned to the generators;
every relation is homogeneous for these degrees.  This constructs
the same algebra for arbitrary integer degrees, including degree zero.

The differential formula is integral.  One way to verify all its
compatibilities without cancelling integers in $k$ is to work first with
formal differential coefficients over a polynomial ring over $\mathbb Z$.
The displayed basis makes that algebra torsion free.  After tensoring with
$\mathbb Q$, the ordinary derivation rule gives the formula for divided
powers; differentiating a basis monomial has integral coefficients and
lies in the original algebra.  Injectivity into the rational algebra
therefore proves the PD derivation identity integrally.  Specializing the
formal coefficients proves it over $k$.  Each calculation involves
finitely many generators and coefficients, so this argument also covers
an infinite basis.  The square of an odd derivation is a derivation.
On an even generator's divided power its square is
\[
d^2(b^{[n]})=b^{[n-2]}(d_Xb)^2+b^{[n-1]}d_X^2b=0
\]
for $n\geq2$; the first term vanishes by the strict odd-square identity.
On $X$ it is zero by $d_X^2=0$, hence it is zero everywhere.
The augmentation ideal is differential stable because $d_X$ takes
generators to linear combinations of generators.

For a chain map $\rho:X\to J$, define its extension by
\[
x^{[f]}\longmapsto\prod_{b\in B}\gamma_{f(b)}(\rho(b)),           \tag{B.5}
\]
where only indices zero and one occur for odd generators.
The multiplication and PD identities make this a graded PD morphism.
Differential compatibility follows on generators from the chain-map
identity and on their products and divided powers from the two
derivation formulas.  Every monomial is obtained by these operations,
which also proves uniqueness.  Restriction and extension are inverse,
establishing \textup{(B.4)}.

To compute $j_X$, decompose the basis of $X^{\otimes n}$ into permutation
orbits indexed by multisets.  An invariant element on one orbit is
determined by the coefficient of its ordered word.  A repeated odd basis
element gives a stabilizing transposition of sign $-1$, so this coefficient
must lie in $k[2]$.  That condition is sufficient: different choices of a
permutation representing the same word change the sign by $-1$, which
acts trivially on $k[2]$.  Without odd repetition all stabilizer signs are
$+1$, and the signed orbit sum has arbitrary coefficient in $k$.
The shuffle product in \textup{(B.5)}, with the generators now tensors,
sends each basis monomial to exactly this signed orbit sum, with
coefficient one.  It therefore embeds the direct summand indexed by
multisets without odd repetition and gives the asserted graded cokernel.
Here TS carries the shuffle divided powers: for a tensor of even
internal degree and fixed tensor length, we sum its tensor powers over
the shuffles of unordered blocks, and extend to even sums by the PD
sum formula.  These finite integral polynomial formulas commute with
specialization of coefficients.  To check divided powers on an arbitrary
even element in the image of $j_X$, express it as a finite sum of signed
orbit monomials and replace its coefficients by independent variables.
Over the resulting polynomial ring over $\mathbb Z$, both algebras are
torsion free, their displayed orbit bases agree, and their PD operations
agree after tensoring with $\mathbb Q$, where they equal $z^n/n!$.
They therefore agree before tensoring, and specializing the coefficients
proves the assertion over $k$, including its $2$-torsion.
The differential formulas agree on generators and their divided powers
and extend by Leibniz.  Thus $j_X$ is a morphism in the category of
strict PD dg triples, with the tensor-positive augmentation ideal in TS.
Naturality follows from the extension property.
The cokernel assertion is about graded modules; no splitting as
cochain complexes is used.
\end{proof}

The coefficient condition in this comparison is essential.  For
$k=\mathbb F_2$ and one generator $e$ in degree $-1$ with $de=0$,
$F_k(ke)$ has zero component in tensor weight two, while
$\mathrm{TS}_k^2(ke)=k(e\otimes e)$.  More explicitly, to the square-zero
PD triple $k\oplus ky$, $|y|=-2$, there are two PD dg maps from
$\mathrm{TS}_k(ke)$ extending the zero map on $ke$: the augmentation, and
the map sending $e^{\otimes2}$ to $y$ and every other positive tensor
length to zero.  Multiplication is respected because the only two
positive lengths summing to two are $1,1$, whose shuffle coefficient is
$2=0$; all higher products map to zero.  A higher divided power of an
even positive tensor has tensor length at least four, so both sides of
PD compatibility are zero.  This proves the failure of the unrestricted
TS freeness assertion at an exact map.  Our correction uses $F_k$ for all
rings and uses $j_X$ to identify it with TS whenever justified.

\subsection{The exact obstruction and the corrected lifting statements}

For complexes $X,K$ we use the product Hom complex
\[
\operatorname{Hom}^r_k(X,K)
=\prod_{i\in\mathbb Z}\operatorname{Hom}_k(X^i,K^{i+r}),\qquad
\partial a=d_Ka-(-1)^r a d_X.                                   \tag{B.6}
\]
Thus unbounded complexes entail products in Hom; the algebras themselves
still contain finite sums of monomials.

\begin{theorem}[B-003: Proved; complete obstruction for a given lift]
\label{B003}
Let $\pi:(R',J',\gamma')\to(R,J,\gamma)$ be a PD dg morphism
surjective on ideals, let $X$ have projective terms, and let
$\rho:X\to J$ be a chain map.  Put $K=\ker(J'\to J)$.
For any graded lift $g:X\to J'$ of $\rho$, the formula
\[
o(\rho,\pi)=[d_{J'}g-gd_X]\in H^1\operatorname{Hom}_k(X,K)       \tag{B.7}
\]
defines a class independent of $g$.
The set of chain lifts is nonempty exactly when this class is zero;
after one lift is fixed, all lifts are obtained uniquely by adding a
degree-zero cocycle of $\operatorname{Hom}_k(X,K)$.
For termwise-free $X$, these chain lifts correspond bijectively to PD dg
lifts of the map $F_k(X)\to R$ induced by $\rho$.
\end{theorem}

\begin{proof}
Termwise projectivity gives $g$.  Applying $\pi$ to its defect gives
$d_J\rho-\rho d_X=0$, so the defect lands in $K$.
Its Hom differential is zero because $d_{J'}^2=d_X^2=0$.
Replacing $g$ by $g+h$, $h\in\operatorname{Hom}^0(X,K)$, changes the
defect by $\partial h$.  This proves independence of the class.
A chain lift has zero defect.  Conversely, an equality
$d_{J'}g-gd_X=\partial h$ gives the chain lift $g-h$.
The difference between two chain lifts lands in $K$ and has zero
differential; adding any such difference preserves both required
identities.  Finally apply the explicit restriction and extension
bijection \textup{(B.4)}.  Equality after composing with $\pi$ is tested
on the generators by uniqueness in \textup{(B.4)}.
\end{proof}

\begin{theorem}[B-004: Proved; exact universal lifting hypotheses]
\label{B004}
Fix a cochain complex $X$ of projective $k$-modules.
\begin{enumerate}
\item Every chain map from $X$ lifts through every degreewise
surjection of complexes precisely when $X$ is contractible.
\item Every chain map from $X$ lifts through every degreewise
surjective quasi-isomorphism precisely when $X$ is K-projective,
meaning that $\operatorname{Hom}_k(X,K)$ is acyclic for every acyclic $K$.
\end{enumerate}
For termwise-free $X$, the same two characterizations hold with chain
maps replaced by PD dg maps out of $F_k(X)$ and with the tested
morphisms required to be surjective, respectively surjective
quasi-isomorphisms, on their PD ideals.  Every bounded-above complex of
projectives meets the hypothesis in the second assertion.
\end{theorem}

\begin{proof}
For every complex $Y$, form
\[
C(Y)^n=Y^n\oplus Y^{n-1},\quad
d_C(a,b)=(d_Ya,a-d_Yb),\quad p(a,b)=a.                           \tag{B.8}
\]
This is a complex, $p:C(Y)\to Y$ is a degreewise-surjective chain map,
and $h_C(a,b)=(b,0)$ has degree $-1$ with $d_Ch_C+h_Cd_C=1$.
Universal lifting applied to $p:C(X)\to X$ and $1_X$ gives a chain
section $s$.  Then $ph_Cs$ contracts $X$.
Conversely, for a contraction $h$ of $X$, a graded lift $g:X\to E$
through $q:E\twoheadrightarrow Y$ produces
\[
\ell=d_Egh+ghd_X:X\longrightarrow E.
\]
The map $\ell=\partial(gh)$ is a cocycle, and
$q\ell=\rho(d_Xh+hd_X)=\rho$.  This proves the first assertion,
including its necessity.

For the second assertion, the kernel $K$ of a surjective
quasi-isomorphism is acyclic, by the long exact cohomology sequence.
K-projectivity therefore kills the obstruction in \textup{(B.7)}.
Conversely let $Y$ be acyclic.  Both $C(Y)$ and $Y$ are acyclic,
so \textup{(B.8)} is a surjective quasi-isomorphism.  Every chain map
$f:X\to Y$ consequently factors as $p\ell$ for a chain map
$\ell:X\to C(Y)$ and is null-homotopic via $ph_C\ell$.
Apply this also to each shift of $Y$.  The cycles of degree $r$ in
\textup{(B.6)} identify, with the stated shift signs, with chain maps
$X\to Y[r]$, and their boundaries identify with null-homotopic maps.
Thus every $H^r\operatorname{Hom}_k(X,Y)$ is zero.

For the PD assertions, sufficiency follows from \textup{(B.4)}.
Necessity uses the square-zero functor
$S:\mathrm{Cplx}(k)\to\mathrm{PDdgAlg}(k)$,
$S(Y)=(k\oplus Y,Y,\gamma_{>1}=0)$, whose identities were checked
in the proof of B-001.  Every map $q:E\to Y$ gives $S(q)$ with exactly
the same map on ideals.  A PD lift restricts to a chain lift, so testing
these actual triples already forces both necessary conditions.

For the final assertion we use the known null-homotopy result
\cite[Lemma 13.19.4]{B_StacksProjective}: a map from a bounded-above
projective complex to an acyclic complex is null-homotopic.  Applied
to each shift of the acyclic complex, with \textup{(B.6)}, it states
exactly that its Hom complex is acyclic.  Equivalently, the descending
construction lifts at degree $n$ the closed remainder
$a^n-(-1)^r h^{n+1}d_X^n$ through
$d:K^{n+r-1}\twoheadrightarrow Z^{n+r}K$.
The construction begins above the bound and uses projectivity at every
degree; it imposes no lower bound and no finite-rank assumption.
\end{proof}

These are necessary-and-sufficient hypotheses for the two stated
universal lifting properties, and B-003 is necessary and sufficient
for each individual lifting problem.  In the ring and complex of B-001,
$K=kv$ in degree $-1$ and \textup{(B.7)} is the class of
$\varepsilon\mapsto v$; its Hom complex has no degree-zero term that
could kill this class.

\begin{proposition}[B-005: Known example; proved square-zero realization]
\label{B005}
Termwise freeness alone remains insufficient for lifting through a
surjective quasi-isomorphism for unbounded complexes.
\end{proposition}

\begin{proof}
The standard example in \cite[p.~124, example (1) and the following
definition]{B_Spaltenstein1988} has $k=\mathbb Z/4$,
$X^n=k$ for every $n\in\mathbb Z$, and $d_X^n=2$.
It is acyclic, since both the kernel and image of multiplication by two
are $2k$.  A contraction would require
$2h^n+2h^{n+1}=1$ in $k$ in every degree, which is impossible.
The map $p:C(X)\to X$ in \textup{(B.8)} is a surjective
quasi-isomorphism with no chain section, because a section would give
the contraction $ph_Cs$.
Applying $S$ gives a morphism of actual square-zero PD dg triples,
surjective and a quasi-isomorphism on ideals, for which the map
$F_k(X)\to S(X)$ induced by $1_X$ has no PD dg lift.
The program checks $\ker(2)=\operatorname{im}(2)$ and all sixteen
possible scalar pairs in the local contraction equation over $k$.
The complex is unbounded; this finite calculation certifies the
degreewise contradiction, not a finite truncation replacement.
\end{proof}

\subsection{Dependence and the resolution comparison}

The source audit removed suppressed \verb|\delete{...}| material and
comments while preserving source lines.  Proposition 3.10 has no label,
and the published text has no later invocation of its number.
The explicit comparison dependencies are
\[
3.15\longrightarrow3.16\longrightarrow3.17,
\qquad \{3.15,3.16\}\longrightarrow4.13.
\]
The last proof also cites Conjecture 4.8; we make no assertion resolving
that conjecture.  Theorem 4.15 uses a separate boundary-lifting induction.
Theorem 3.11 and the construction in 3.13 use explicit adjunctions.
These inspections identify no use of Proposition 3.10 by another
numbered result.  Their recorded source locators are in
\texttt{B\_source\_audit.json}.

We now prove the actual comparison needed to assess the reach of D-031.
For this paragraph we use the Tate convention of
\cite[Definitions 23.6.1 and 23.6.5]{B_StacksTate}, in cochain degrees:
a nonpositive strictly graded commutative dg algebra has divided powers
on its strictly negative even part, compatible with $d$.
This convention permits $d(P^{-1})\subseteq P^0$.

\begin{theorem}[B-006: Known comparison; explicit derivation]
\label{B006}
Let $A$ be any commutative ring and $I\subseteq A$ any ideal.
Let $P\to A/I$ be a Tate resolution constructed as follows.
Start with $P_0=A$.  For each $n\geq0$, adjoin a set of free generators
$F_{n+1}$ in cochain degree $-(n+1)$, exterior for odd degree and with
divided powers for even degree, whose differentials are prescribed
cycles $z_b\in P_n^{-n}$; for $n=0$ these lie in $I$.
The generators kill the required cohomology and
$P=\bigcup_nP_n\to A/I$ is a quasi-isomorphism.

Let $Q$ be a nonpositive Tate PD dg $A$-algebra and
$q:Q\to A/I$ an augmented quasi-isomorphism.
There is an augmentation-preserving PD dg $A$-algebra morphism
$\theta:P\to Q$.  For two such free resolutions the two comparison
maps compose to maps chain homotopic to the identities as $A$-module
maps.
\end{theorem}

\begin{proof}
This is the form of \cite[Lemma 23.6.11, Tag 09PP]{B_StacksTate}
obtained by the exact reindexing $P_m=P^{-m}$, $Q_m=Q^{-m}$.
The reindexing leaves the differential map itself unchanged,
turns degree $+1$ into homological degree $-1$, and preserves all
Koszul signs because $(-i)(-j)\equiv ij\pmod2$.
We give the extension argument to specify its application and correct
the intermediate-map convention.

The structure map $\theta_0:A\to Q^0$ is the required map at stage zero.
For a first-stage generator $b$, the element $\theta_0(z_b)$ has zero
image in $A/I$.  Since $q$ induces $H^0(Q)\cong A/I$ and $Q^{>0}=0$,
it belongs to $d_Q(Q^{-1})$.  Choose $y_b\in Q^{-1}$ with
$d_Qy_b=\theta_0(z_b)$.
At stage $n\geq1$, a previously constructed chain map $\theta_n$
sends $z_b$ to a cycle in $Q^{-n}$.  The equality $H^{-n}(Q)=0$
gives $y_b\in Q^{-(n+1)}$ with $d_Qy_b=\theta_n(z_b)$.
These choices define a linear map on the free module $F_{n+1}$.

Extend $b\mapsto y_b$ and $\theta_n$ as a graded PD algebra map
$\theta_{n+1}:P_{n+1}\to Q$, by the explicit ordered-monomial formula
\textup{(B.5)} over $P_n$.  The differential identity holds on $P_n$
and every new generator.  It propagates to products by Leibniz and to
each even generator's divided powers by
\[
d_Q\gamma_r(y_b)
=\gamma_{r-1}(y_b)\theta_n(z_b)
=\theta_{n+1}(b^{[r-1]}z_b).
\]
This proves it on a basis, hence on all of $P_{n+1}$.
The maps agree on preceding stages and on the augmentation, so their
union is the required map $\theta:P\to Q$.
The auxiliary map from the isolated zero-differential free algebra on
the new generators is a \emph{graded} PD map; the chain-map assertion
has just been proved for the full extension with $db=z_b$.

For comparison maps $\theta:P\to Q$ and $\tau:Q\to P$, the map
$f=\tau\theta-1_P$ lands in $L=\ker(P\to A/I)$.
This kernel is acyclic and $P$ is bounded above with free terms.
By B-004, $\operatorname{Hom}_A(P,L)$ is acyclic.
Thus $f=d_Ph+hd_P$ for an $A$-linear degree-$-1$ map $h:P\to L$.
The same argument gives the other homotopy.
These are homotopies of module complexes; no multiplicative homotopy
is asserted or used.
\end{proof}

For D-031, the map $k\oplus E\to k\oplus M$ has source cohomology
$k$ in degree zero and target cohomology $k$ in degrees zero and $-2$.
Consequently it cannot supply the augmented quasi-isomorphism in B-006
or in the target-resolution hypothesis of Theorem 3.15.
The identities $d_Qy_b=\theta_n(z_b)$ in the proof specify the precise
additional surjections that make the comparison possible.

\begin{proposition}[B-007: Refuted convention; proved correction]
\label{B007}
The strictly negative part of the Tate resolution used in the paper's
Theorem 3.13 need not be a differential ideal in the sense of its
\S2.2.  The Tate convention used in B-006 gives the required
resolution and comparison with all coefficient rings retained.
\end{proposition}

\begin{proof}
Take $A=\mathbb Z$ and $I=(2)$.  The first step has
$P_1=\bigwedge_A(e)$, $|e|=-1$, $de=2$.
Thus $e\in P_1^{<0}$ but $de\notin P_1^{<0}$.
Every later stage contains this same equality, so $P^{<0}$ is not a
differential ideal.  The nonpositive dg algebra $P_1$ is already a
resolution of $A/I$, because multiplication by two in $A$ is injective.
Its positive-index divided powers on the negative even part are defined
on the zero module, so it is a Tate PD dg algebra exactly as in B-006.
The augmentation $P_1\to\mathbb Z/2$ is the explicit comparison to
the quotient in degree zero.

For any nonpositive strict PD triple $(R,J,\gamma)$ with $R^{<0}\subseteq J$,
restriction of $\gamma$ to $R^{<0}_{\rm even}$ defines a Tate object;
the underlying morphism is the identity on $R$ and the operations are
literally restricted.  Its derivative identity is unchanged.  The above
resolution proves that this forgetful passage does not supply a
differential ideal equal to $P^{<0}$ in the reverse direction.
B-006 states the comparison entirely in the resulting Tate category
and proves it there.  Together with B-002 this identifies the two
repairs needed before treating the paper's resolution construction as
a uniform literal assertion over arbitrary coefficient rings.
\end{proof}

\begin{proposition}[B-008: Proved; the cohomology algebra comparison]
\label{B008}
For the comparison maps of B-006, put
$\alpha(D)=\theta D\tau$ and $\alpha'(D')=\tau D'\theta$ on the
endomorphism Hom complexes.  Then $H^*(\alpha)$ and $H^*(\alpha')$
are mutually inverse isomorphisms of graded $A$-algebras.
\end{proposition}

\begin{proof}
The two maps are chain maps by the Hom differential formula.
Write $\tau\theta-1_P=\partial h$, $|h|=-1$.
For cocycles $D_1,D_2$ of degrees $a,b$ in $\operatorname{Hom}_A(P,P)$,
the discrepancy between multiplication and $\alpha$ is
\[
\begin{split}
\alpha(D_1)\alpha(D_2)-\alpha(D_1D_2)
 &=\theta D_1(\tau\theta-1_P)D_2\tau\\
 &=(-1)^a\partial(\theta D_1hD_2\tau).
\end{split}
\]
The Hom differential satisfies its signed Leibniz identity by direct
substitution in \textup{(B.6)}, proving the second equality.
Thus $H^*(\alpha)$ is multiplicative.  The identity
$[\tau\theta]=[1_P]$ gives
$[\alpha'\alpha(D)]=[\tau\theta][D][\tau\theta]=[D]$,
and the other homotopy gives the reverse composite and preservation
of the unit.  This proves the stated isomorphisms with their products,
which is the comparison underlying Corollary 3.17.
\end{proof}

Accordingly D-031 is confirmed, with B-003 and B-004 as exact corrected
lifting statements.  Its failure reaches neither the independently
proved Tate comparison B-006 nor B-008.  This conclusion does not certify
the paper's other claims: B-002 and B-007 record the precise additional
convention issues encountered while checking the hypotheses used here.
